% 第5章图：由消元次序(1,2,5,3,4)构造树分解。边{1,2},{1,4},{2,3},{3,4},{3,5},{4,5}。
% 袋：B1={1,2,4}（补边{2,4}），B2={2,3,4}，B5={3,4,5}，B3={3,4}，B4={4}；父袋 p(1)=2,p(2)=3,p(5)=3,p(3)=4。
\begin{tikzpicture}[x=1mm,y=1mm,>=Stealth,
  v/.style={circle,draw=FigureStroke,fill=FigureBlueFill,line width=.7pt,minimum size=4.8mm,inner sep=0pt,font=\figurefont\fontsize{7.5}{9}\selectfont},
  cur/.style={v,fill=FigureYellowFill,draw=FigureYellow},
  e/.style={FigureStroke,line width=.7pt},
  fe/.style={FigureGreen,line width=.9pt,dashed},
  bag/.style={rectangle,rounded corners=1.8pt,draw=FigureStroke,fill=FigureBlueFill,line width=.7pt,minimum height=6mm,inner sep=1.5pt},
  sc/.style={font=\figurefont\footnotesize,text=FigureInk},
  note/.style={font=\FigureNoteFont\fontsize{7.5}{9}\selectfont,text=FigureMuted}]
\def\S{23}
% step 1
\begin{scope}[yshift=40mm]
 \node[cur] (a1) at (0,12) {1}; \node[v] (a2) at (10,12) {2}; \node[v] (a3) at (10,2) {3}; \node[v] (a4) at (0,2) {4}; \node[v] (a5) at (5,-6) {5};
 \draw[e] (a1)--(a2) (a2)--(a3) (a3)--(a4) (a1)--(a4) (a3)--(a5) (a4)--(a5);
 \draw[fe] (a2)--(a4);
 \node[sc] at (5,-13) {$\{1,2,4\}$};
 \node[note] at (5,19) {处理$1$};
\end{scope}
\begin{scope}[xshift=\S mm,yshift=40mm]
 \node[cur] (a2) at (10,12) {2}; \node[v] (a3) at (10,2) {3}; \node[v] (a4) at (0,2) {4}; \node[v] (a5) at (5,-6) {5};
 \draw[e] (a2)--(a3) (a3)--(a4) (a3)--(a5) (a4)--(a5);
 \draw[FigureGreen,line width=.9pt] (a2)--(a4);
 \node[sc] at (5,-13) {$\{2,3,4\}$};
 \node[note] at (5,19) {处理$2$};
\end{scope}
\begin{scope}[xshift=2*\S mm,yshift=40mm]
 \node[v] (a3) at (10,2) {3}; \node[v] (a4) at (0,2) {4}; \node[cur] (a5) at (5,-6) {5};
 \draw[e] (a3)--(a4) (a3)--(a5) (a4)--(a5);
 \node[sc] at (5,-13) {$\{3,4,5\}$};
 \node[note] at (5,19) {处理$5$};
\end{scope}
\begin{scope}[xshift=3*\S mm,yshift=40mm]
 \node[cur] (a3) at (10,2) {3}; \node[v] (a4) at (0,2) {4};
 \draw[e] (a3)--(a4);
 \node[sc] at (5,-13) {$\{3,4\}$};
 \node[note] at (5,19) {处理$3$};
\end{scope}
\begin{scope}[xshift=4*\S mm,yshift=40mm]
 \node[cur] (a4) at (0,2) {4};
 \node[sc] at (0,-13) {$\{4\}$};
 \node[note] at (0,19) {处理$4$};
\end{scope}
% bag tree
\begin{scope}[yshift=4mm]
 \node[bag] (B1) at (2,8) {$B_1=\{1,2,4\}$};
 \node[bag] (B2) at (32,8) {$B_2=\{2,3,4\}$};
 \node[bag] (B3) at (62,8) {$B_3=\{3,4\}$};
 \node[bag] (B4) at (90,8) {$B_4=\{4\}$};
 \node[bag] (B5) at (62,-8) {$B_5=\{3,4,5\}$};
 \draw[->,FigureStroke,line width=.8pt] (B1)--(B2) node[midway,above,note]{$\{2,4\}$};
 \draw[->,FigureStroke,line width=.8pt] (B2)--(B3) node[midway,above,note]{$\{3,4\}$};
 \draw[->,FigureStroke,line width=.8pt] (B3)--(B4) node[midway,above,note]{$\{4\}$};
 \draw[->,FigureStroke,line width=.8pt] (B5)--(B3) node[midway,right,note]{$\{3,4\}$};
 \node[note,anchor=west] at (B4.south west) [yshift=-3mm] {根};
\end{scope}
\end{tikzpicture}
